\section{The Capstone: The Model Run End-to-End (P34, P36)}
\label{sec:capstone}

Every result above used the substrate best suited to it. The honest question is whether
it all falls out of one instantiation of the present model of Section~\ref{sec:model}.

\subsection{The capstone run (P34)}

We run that one model, a growing random hyperbolic mesh with the ternary signed-majority
rule, and read every key emergent structure off the same mesh. \textbf{Where we landed}
(one mesh of $1500$ vibes):
\begin{itemize}
\item the mesh has mean degree $10.6$, Lorentz anisotropy $0.070$ (no preferred frame),
and exponential reach,
\item the dynamics keeps tones strictly ternary and converges to a stable structured
state (the flip fraction falls to zero), a non-trivial fixed point in the manner of a
Hopfield network, the basis for persistent matter and selves,
\item the emergent Hamiltonian (the graph Laplacian) is bounded below (minimum
eigenvalue $0.007 \ge 0$) and local (range $1$),
\item the arrow of time holds, relations only accumulate as the mesh grows.
\end{itemize}
So the model runs as one system. From a single instantiation come a definite geometry
with no preferred frame, stable structured states, the local bounded-below Hamiltonian,
and the arrow. The matter, force, and gravity sectors are read off the same mesh by
their operators. The natural next integration is to evolve all of them by the one
microscopic rule in a single simulation.

\subsection{The model written at a glance (P36)}

The present model is written and run by a small constructor, so its description and
its implementation are one object. With no options it is the present model, and a
one-word change expresses a variant for comparison:
\begin{verbatim}
  const model = vibe().size(1500).seed(1)    // the present model
  console.log(model.describe())              // print it at a glance
  const world = model.build().run(40)        // build the mesh, run 40 beats
  world.read()                               // emergent structures off the mesh
\end{verbatim}
Swapping \texttt{vibe().mesh('lattice')} raises the Lorentz anisotropy from $0.06$ to
$1.0$ (a strong preferred frame), making plain why the random hyperbolic mesh is the
one chosen: the lattice variant breaks Lorentz invariance and the hyperbolic one
does not. This is the model definition the rest of the paper rests on, legible and
executable at once.
